\section{Разложение функции. Ряд Тейлора}\label{sec-power-series-04}

\begin{theorem}[Коэффициенты разложения]\label{thm-power-coefficients}

Если \(S(z)=\sum_{n=0}^\infty c_n(z-z_0)^n\) имеет положительный радиус сходимости, то

\begin{equation}
c_n=\frac{S^{(n)}(z_0)}{n!},\qquad n\ge0.
\label{eq-taylor-coefficients}
\end{equation}

\end{theorem}

\begin{proof}

\(S(z_0)=c_0\). В формуле \eqref{eq-power-higher-derivatives} подставляем \(z=z_0\): единственный ненулевой член соответствует \(n=k\), поэтому \(S^{(k)}(z_0)=k!c_k\).

\end{proof}

\begin{cor}[Единственность разложения]\label{cor-power-uniqueness}

Если на некотором круге \(|z-z_0|<r\), \(r>0\), \[
\sum_{n=0}^\infty a_n(z-z_0)^n
=\sum_{n=0}^\infty b_n(z-z_0)^n,
\] то \(a_n=b_n\) для всех \(n\ge0\).

\end{cor}

Это следует из формулы \eqref{eq-taylor-coefficients}. Для ряда с центром в нуле: если сумма чётна, все коэффициенты нечётных степеней равны нулю; если нечётна, все коэффициенты чётных степеней равны нулю.

При другом центре нужно рассматривать симметрию относительно \(z_0\).

\begin{definition}[Ряд Тейлора и ряд Маклорена]\label{def-taylor-series}

Функции \(f\), имеющей производные всех порядков в точке \(z_0\), сопоставляют ряд

\begin{equation}
\sum_{n=0}^\infty\frac{f^{(n)}(z_0)}{n!}(z-z_0)^n.
\label{eq-taylor-series}
\end{equation}

Он называется рядом Тейлора функции в точке \(z_0\), а при \(z_0=0\) --- рядом Маклорена.

\end{definition}

Для вещественной бесконечно дифференцируемой функции радиус этого ряда и совпадение его суммы с \(f\) требуют отдельного исследования. Любой сходящийся степенной ряд является рядом Тейлора своей суммы.

Для комплексной функции, голоморфной в окрестности, разложение в ряд Тейлора там существует. Предостережение о несовпадении суммы с функцией относится к общему вещественному \(C^\infty\)-случаю.

\begin{definition}[Аналитичность]\label{def-analytic}

Функция называется аналитической (регулярной) в точке \(z_0\), если в некоторой её окрестности допускает представление сходящимся степенным рядом с центром \(z_0\).

\end{definition}

\subsection{Достаточные условия разложения вещественной функции}\label{sec-power-series-04-1}

\begin{theorem}[Оценки производных]\label{thm-taylor-sufficient}

Пусть \(f:\mathbb R\to\mathbb R\) бесконечно дифференцируема в окрестности отрезка \(|x-x_0|\le r\), \(r>0\). Каждого из следующих условий достаточно для совпадения \(f\) со своим рядом Тейлора на этом отрезке:

\begin{enumerate}
\def\labelenumi{\arabic{enumi}.}
\tightlist
\item
  \(0<r<1\) и существует \textbf{одна} константа \(M>0\) такая, что \(|f^{(n)}(x)|\le M n!\) для всех \(n\ge0\) и всех \(|x-x_0|\le r\).
\item
  Существует \textbf{одна} константа \(M>0\) такая, что \(|f^{(n)}(x)|\le M\) для всех \(n\ge0\) и всех \(|x-x_0|\le r\).
\end{enumerate}

\end{theorem}

\begin{proof}

По формуле Тейлора с остатком Лагранжа при некотором \(c\) между \(x\) и \(x_0\), \[
f(x)-\sum_{k=0}^{n-1}\frac{f^{(k)}(x_0)}{k!}(x-x_0)^k
=\frac{f^{(n)}(c)}{n!}(x-x_0)^n.
\] В первом случае модуль остатка не превосходит \(Mr^n\to0\), во втором --- \(Mr^n/n!\to0\). Оценки равномерны по \(|x-x_0|\le r\). Поэтому в обоих случаях \[
f(x)=\sum_{k=0}^\infty\frac{f^{(k)}(x_0)}{k!}(x-x_0)^k.
\]

\end{proof}
